%! Unit Launch: Study Design — Unit 1, Lesson 1.0 — Guided Notes & Practice (Answer Key)
% THE in-class document, in the MAIN IDEAS / NOTES shape (LESSON_SHAPE.md §1, ported from AP
% Statistics 2026-09-12): the vocabulary box, the hook, then ONE two-column guidednotes table.
% Each row is one idea — a label on the left; on the right one or two COMPLETE printed
% sentences, ONE large pre-drawn display the student reads or names, then a bold prompt and a
% two-across grid of numbered problems with 1.2–1.8cm of writing room. Problems run 1..17.
% DENSITY: a \blank{} only where a word IS the answer (the four cells of the classify table
% in row 3, problem 5 — the one \wordbank strip sits above it); no mid-sentence blanks; every
% display is pre-drawn; the page belongs to the student's pen.
%   I DO   : the teacher reads the definition, marks up the display, works the first problem
%            of each grid (1, 3, 6, 11).
%   WE DO  : the class works the rest of each grid, students holding the pen; the last row,
%            Guided Practice (14–17), is worked entirely together on the Student Life Survey
%            instead of the basketball roster.
%   YOU DO : nothing on this page — ../homework/ IS the individual practice, started in the
%            last ten minutes of the period. No independent set, no reflection box, no
%            objectives box (the cover carries the targets).
% THE TRAP is problem 9 (row 3, "Two Types"): Jersey # is classified LAST, by hand vote, and
% NOT settled — the digit reflex says quantitative. THE CRUX is problem 10 (row 4, "The Digit
% Test"): the blog's mean of 14 is arithmetically correct and describes nobody — the hook vote
% is re-taken there (PS.DC.1c). Problem 15 is the crux transferred to the survey's Grade
% column (63/6 = 10.5); problem 13 is the counterweight (Height is digits and averages fine).
%
% THE DATA (all verified in Python 2026-09-14)
%   Riverbend varsity roster, 6 players: Jersey 12,7,23,4,17,21 -> 84/6 = 14 (nobody wears 14);
%     Points 14,9,6,19,11,13 -> 72/6 = 12; Height 64,71,75,68,70,72 -> 420/6 = 70; captains 3/6 = 50%
%   Grade levels 9,10,11,12 -> 42/4 = 10.5; ZIPs 22030, 22201 -> 22115.5; siblings 2,0,3,1,2,4 -> 12/6 = 2
%   Guided Practice, Student Life Survey (6 students): Commute 20,8,35,15,5,25 -> 108/6 = 18 min;
%     Grade 10,11,9,12,10,11 -> 63/6 = 10.5
%
% PAGE LOCKSTEP with notes_key/: the key is GENERATED from this file by templates/lesson/mkkey.py
% (spec: templates/lesson/examples/lesson00_notes_key.json) — every \blank{} becomes an \ans{},
% every \vterm{} a \vtermans{}{}, every \par\writelines{n} n \ansline{} calls, and the answer
% argument of every \pcell, \writespace and \labelbox is filled. Answer spaces are fixed
% heights, so the two files cannot drift. KEEP EVERY \pcell ON ONE SOURCE LINE — mkkey matches
% line by line.
% TABLE MECHANICS: inside a cell, `\\` ENDS THE ROW — break lines with \par. A row cannot break
% across pages: the figure gets its own sub-row (`\\` then `& ...`) and EACH GRID ROW is its own
% sub-row — close the probgrid, `\\ &`, reopen as probgrid* (no top rule).
\documentclass[12pt]{article}
\usepackage{saar-article}
\usepackage{saar-key}   % pulls in -boxes; do NOT also load -boxes
\newcommand{\TallMath}[1]{$\displaystyle #1\rule[-1.4em]{0pt}{3.2em}$}
% Word bank strip — ONLY above a table the student fills (row 3, problem 5). Defined per
% document; the key inherits it and prints the same strip, so heights match.
\newcommand{\wordbank}[1]{\par\vspace{2pt}\noindent\fcolorbox{goldacc}{goldbg}{\parbox{\dimexpr\linewidth-2\fboxsep-2\fboxrule\relax}{\small\textbf{Word bank:} #1}}\par\vspace{4pt}}
% Vocabulary rows: a FIXED-HEIGHT row carrying the term label and OPEN writing space —
% no underline beside the term and no rule line (user correction 2026-09-06: the black
% answer line crowded the row; students write beside and below the term). The key fills
% exactly the same height, so the box cannot drift blank vs. keyed. saar-article's
% \termblank adds an inline blank and a rule line, so the pair is defined here (the key is
% generated from this file and inherits both). 2.3cm per row gives students three
% handwritten lines of room; keep key definitions to two lines.
\newcommand{\termrowheight}{2.3cm}
\newlength{\termrowinset}\setlength{\termrowinset}{7pt}
\newcommand{\vterm}[1]{%
  \par\noindent\begin{minipage}[t][\termrowheight][t]{\linewidth}%
    \vspace*{\termrowinset}%
    \textbf{\textcolor{cerulean}{#1:}}%
  \end{minipage}\par}
\newcommand{\vtermans}[2]{%
  \par\noindent\begin{minipage}[t][\termrowheight][t]{\linewidth}%
    \vspace*{\termrowinset}%
    \textbf{\textcolor{cerulean}{#1:}}\quad\ans{#2}%
  \end{minipage}\par}

\begin{document}

\pageheader{Unit 1, Lesson 1.0}{Guided Notes \& Practice --- Answer Key}
% NAMESTRIP: no \namedateperiod here — mirrors the blank. NO teachernote here —
% teacher prose lives in the lesson plan.

\vspace{0.15in}

\begin{vocabbox}
\small
Fill in each term \textbf{as we name it} in the notes below.
\vtermans{Data}{Numbers or labels collected about individuals, together with the context --- who, what, and the units --- that gives them meaning.}
\vtermans{Individual}{One person, animal, or object that data are collected about --- one row of a data table.}
\vtermans{Variable}{A characteristic recorded for each individual --- one column of a data table. The name column is not one.}
\vtermans{Categorical variable}{A variable that records a label or group name (position, lunch, yes/no). Summarize it with counts and percents --- a mean means nothing.}
\vtermans{Quantitative variable}{A variable that records a measured amount with units (points, inches, minutes). Summarize it with a mean.}
\end{vocabbox}

\boxguard[12]
\begin{hookbox}
\small
\textit{``The Riverbend roster's mean jersey number is $14$. That is low for a varsity squad
--- this is a young, inexperienced team.''} --- a sports blog. The arithmetic is \textbf{correct}.

\vspace{3pt}
\textbf{Is the conclusion correct?} Circle one: \quad yes \qquad no \qquad can't tell
\quad --- and say why you think so. We vote again at problem~10.
\par\ansline{Answers vary --- most of the room says yes or can't tell; a few say ``jersey numbers don't mean anything.''}
\ansline{(Settled at problem 10: no --- a jersey number is a label, so its mean describes nobody.)}
\end{hookbox}

\small
\begin{guidednotes}
% ── ROW 1: data are numbers with a context ───────────────────────────────────
\mainidea[Numbers with a]{Data in Context} &
\textbf{Data} are numbers or labels together with the context that gives them meaning ---
\textbf{who} was measured, \textbf{what} was measured, and in what \textbf{units}. A bare
number answers nothing. \textbf{A number is not data until it comes with all three.}
\\
& {\centering
\begin{tikzpicture}[scale=0.9]
  % the bare number
  \fill[goldbg, rounded corners=5pt] (0,0) rectangle (2.0,2.5);
  \draw[goldacc, thick, rounded corners=5pt] (0,0) rectangle (2.0,2.5);
  \node[font=\Huge\bfseries] at (1.0,1.5) {12};
  \node[font=\tiny\itshape, align=center] at (1.0,0.45) {a number\\by itself};
  \draw[->, very thick, charcoal] (2.15,1.25) -- (3.05,1.25);
  \node[font=\tiny\bfseries, above, align=center] at (2.6,1.3) {add the\\context};
  % A: who
  \fill[frost, rounded corners=4pt] (3.2,1.4) rectangle (12.4,2.5);
  \draw[cerulean, thick, rounded corners=4pt] (3.2,1.4) rectangle (12.4,2.5);
  \node[font=\footnotesize] at (7.8,1.95) {``The \textbf{six players on the Riverbend varsity roster}};
  \node[font=\scriptsize\bfseries\color{cerulean}, fill=white, inner sep=1.5pt] at (3.55,2.5) {A};
  % B: what
  \fill[goldbg, rounded corners=4pt] (3.2,0) rectangle (7.6,1.1);
  \draw[goldacc, thick, rounded corners=4pt] (3.2,0) rectangle (7.6,1.1);
  \node[font=\footnotesize] at (5.4,0.55) {scored a \textbf{mean of $12$}};
  \node[font=\scriptsize\bfseries\color{goldacc}, fill=white, inner sep=1.5pt] at (3.55,1.1) {B};
  % C: units
  \fill[greenbg, rounded corners=4pt] (7.8,0) rectangle (12.4,1.1);
  \draw[greenacc, thick, rounded corners=4pt] (7.8,0) rectangle (12.4,1.1);
  \node[font=\footnotesize] at (10.1,0.55) {\textbf{points} per player.''};
  \node[font=\scriptsize\bfseries\color{greenacc}, fill=white, inner sep=1.5pt] at (8.15,1.1) {C};
\end{tikzpicture}\par\vspace{5pt}
A is the \labelbox{1.5cm}{who}\qquad B is the \labelbox{1.5cm}{what}\qquad C is the \labelbox{1.5cm}{units}\par}
\\
& \notesprompt{Turn a bare number into data --- name the who, the what, and the units.}
\begin{probgrid}{|Y|Y|}
\pcell{1}{Warm-up item 3 gave you $14$. Write it as data: a sentence with a who, a what, and units.}{1.6cm}{The five students surveyed spent a mean of 14 minutes on homework last night.} & \pcell{2}{The blog's $14$ has a who (six players) and a what (jersey number). What are its \emph{units}? Is $14$ an amount of anything?}{1.6cm}{None --- 14 is not measured in anything. It is not an amount, so the 14 is not really data.} \\ \hline
\end{probgrid}
\\ \hline
% ── ROW 2: individuals and variables ─────────────────────────────────────────
\mainidea[Reading a table]{Individual \& Variable} &
An \textbf{individual} is one person, animal, or object that data are collected about ---
one \textbf{row} of a data table. A \textbf{variable} is a characteristic recorded for each
individual --- one \textbf{column}. \textbf{The name column only tells the rows apart; it is
not a variable we study.}
\\
& {\centering
\begin{tikzpicture}[scale=1.0]
  % the roster the blog was reading; Nia's row and the Points column are shaded
  \fill[frost] (0,-0.9) rectangle (12.2,-0.45);
  \fill[goldbg] (5.9,-3.15) rectangle (7.3,0);
  \draw[cerulean, thick] (0,-0.9) rectangle (12.2,-0.45);
  \draw[goldacc, thick] (5.9,-3.15) rectangle (7.3,0);
  \draw[thick] (0,-0.45) -- (12.2,-0.45);
  \draw[thick] (0,-3.15) -- (12.2,-3.15);
  \foreach \r/\pa/\pb/\pc/\pd/\pe/\pf in {
      0/\textbf{Player}/\textbf{Jersey \#}/\textbf{Position}/\textbf{Points}/\textbf{Height (in)}/\textbf{Captain},
      1/Nia/12/Guard/14/64/Yes,
      2/Ruben/7/Forward/9/71/No,
      3/Sasha/23/Center/6/75/No,
      4/Marcus/4/Guard/19/68/Yes,
      5/Talia/17/Forward/11/70/Yes,
      6/Devon/21/Guard/13/72/No} {
    \node[anchor=west, font=\footnotesize] at (0.1,{-0.45*\r-0.225}) {\pa};
    \node[font=\footnotesize] at (2.6,{-0.45*\r-0.225}) {\pb};
    \node[anchor=west, font=\footnotesize] at (3.7,{-0.45*\r-0.225}) {\pc};
    \node[font=\footnotesize] at (6.6,{-0.45*\r-0.225}) {\pd};
    \node[font=\footnotesize] at (8.9,{-0.45*\r-0.225}) {\pe};
    \node[font=\footnotesize] at (11.2,{-0.45*\r-0.225}) {\pf};
  }
\end{tikzpicture}\par\vspace{5pt}
The shaded row is one \labelbox{2.0cm}{individual}\qquad the shaded column is one \labelbox{2.0cm}{variable}\par}
\\
& \notesprompt{Read the roster. Count, and say what you counted.}
\begin{probgrid}{|Y|Y|}
\pcell{3}{How many individuals does the roster describe? How many variables does it record? (Is ``Player'' one of them?)}{1.6cm}{6 individuals (the players); 5 variables. No --- Player only names the rows.} & \pcell{4}{The athletic director adds a column, \emph{Free throws made}, and a seventh player, Jo. Individuals now? Variables now?}{1.6cm}{7 individuals; 6 variables. Player still is not one.} \\ \hline
\end{probgrid}
\\ \hline
% ── ROW 3: two types of variables — the TRAP is problem 9 ────────────────────
\mainidea[Categorical or quantitative]{Two Types} &
A \textbf{quantitative variable} records a measured \textbf{amount}, with units --- summarize
it with a \textbf{mean}. A \textbf{categorical variable} records a \textbf{label} or group
name --- summarize it with \textbf{counts and percents}. \textbf{The type decides what you are
allowed to do with the column.}
\\
& {\centering
\begin{tikzpicture}[scale=1.0]
  % left panel: quantitative
  \fill[goldbg, rounded corners=5pt] (0,0) rectangle (5.9,2.6);
  \draw[goldacc, thick, rounded corners=5pt] (0,0) rectangle (5.9,2.6);
  \node[font=\footnotesize\bfseries\color{goldacc}] at (2.95,2.25) {QUANTITATIVE --- an amount};
  \node[font=\scriptsize] at (2.95,1.75) {Points: $14,\ 9,\ 6,\ 19,\ 11,\ 13$};
  \node[font=\small\bfseries] at (2.95,1.05) {mean $= \dfrac{72}{6} = 12$ points};
  \node[font=\tiny\itshape] at (2.95,0.3) {you can add them --- so a mean means something};
  % right panel: categorical
  \fill[frost, rounded corners=5pt] (6.3,0) rectangle (12.2,2.6);
  \draw[cerulean, thick, rounded corners=5pt] (6.3,0) rectangle (12.2,2.6);
  \node[font=\footnotesize\bfseries\color{cerulean}] at (9.25,2.25) {CATEGORICAL --- a label};
  \node[font=\scriptsize] at (9.25,1.75) {Position: Guard, Forward, Center, \ldots};
  \node[font=\footnotesize\bfseries] at (9.25,1.05) {$3$ Guards, $2$ Forwards, $1$ Center};
  \node[font=\tiny\itshape] at (9.25,0.3) {you cannot add labels --- so count them instead};
\end{tikzpicture}\par\vspace{5pt}
Height (in) is \labelbox{2.8cm}{quantitative}\qquad Captain is \labelbox{2.8cm}{categorical}\par}
\\
& \textbf{5.} Classify each column of the roster. \emph{(Jersey \# is problem~9 --- leave it for last.)}
\wordbank{categorical \;\textbullet\; quantitative}
\footnotesize\renewcommand{\arraystretch}{1.15}
\begin{tabularx}{\linewidth}{@{} X p{3.4cm} X p{3.4cm} @{}}
\toprule
\textbf{Variable} & \textbf{Type} & \textbf{Variable} & \textbf{Type} \\
\midrule
Position & \ans{categorical} & Points      & \ans{quantitative} \\
Captain  & \ans{categorical} & Height (in) & \ans{quantitative} \\
\bottomrule
\end{tabularx}\small
\\
& \notesprompt{Summarize each kind the way its type allows --- and say how you decided, not ``it's words'' or ``it's numbers.''}
\begin{probgrid}{|Y|Y|}
\pcell{6}{A mean is not available for Captain. Summarize it the other way: what percent of the six players are captains?}{1.4cm}{$3/6 = 0.50 = 50\%$ of the six players are captains.} & \pcell{7}{Find the mean of Points, then say it as data --- who, what, units.}{1.6cm}{$72/6 = 12$. The six players on the roster scored a mean of 12 points per player.} \\ \hline
\end{probgrid}
\\
& \begin{probgrid}*{|Y|Y|}
\pcell{8}{A survey records each student's number of siblings: $2, 0, 3, 1, 2, 4$. Which type? How do you know?}{1.4cm}{Quantitative --- a count, an amount. Mean: 2 siblings.} & \pcell{9}{\textbf{Jersey \#. Vote first:} categorical or quantitative? Record the class split, then your reason. \textbf{We settle it at problem~10.}}{1.6cm}{Split recorded (answers vary). Settled at 10: categorical --- a label written with digits.} \\ \hline
\end{probgrid}
\\ \hline
% ── ROW 4: THE CRUX — the digit test settles the hook ────────────────────────
\mainidea[Settle the vote:]{The Digit Test} &
The blog's arithmetic is right: $\dfrac{12+7+23+4+17+21}{6} = \dfrac{84}{6} = 14$.
\textbf{Now find the player who wears $14$ on average.} There is none --- the mean of a
column of labels describes nobody and measures nothing.
\\
& {\centering
\begin{tikzpicture}[scale=1.0]
  % the six jersey numbers on a number line, and the mean that lands on nobody
  \draw[->, thick] (-0.3,0) -- (12.2,0);
  \foreach \p in {0,5,...,25} { \draw ({\p*0.44},-0.1) -- ({\p*0.44},0.1); \node[below, font=\tiny] at ({\p*0.44},-0.12) {\p}; }
  % labels alternate high/low so the two closest jerseys (21 and 23) do not collide
  \foreach \p/\who/\hi in {4/Marcus/0.16, 7/Ruben/0.62, 12/Nia/0.16, 17/Talia/0.62, 21/Devon/0.16, 23/Sasha/0.62} {
    \fill[cerulean] ({\p*0.44},0) circle (0.11);
    \node[above, font=\tiny\bfseries, align=center] at ({\p*0.44},\hi) {\who\\\#\p}; }
  \draw[redacc, very thick] ({14*0.44},-0.35) -- ({14*0.44},1.45);
  \node[font=\scriptsize\bfseries\color{redacc}, fill=white, inner sep=1.5pt] at ({14*0.44},1.68) {mean $= 14$};
  \node[font=\tiny\itshape\color{redacc}, anchor=west] at (7.4,1.68) {nobody wears it --- and it is not an amount};
\end{tikzpicture}\par}
\\
& \fcolorbox{cerulean}{frost}{\parbox{\dimexpr\linewidth-2\fboxsep-2\fboxrule\relax}{\centering
\textbf{The test is not} ``does it look like a number?'' \textbf{The test is:} \emph{does arithmetic on this variable mean anything?}}}
\par\vspace{3pt}
\textbf{In your own words:} why is a correct mean of the jersey numbers still a worthless number?
\writespace{1.4cm}{Jersey numbers are labels, not amounts. Adding labels means nothing, so their mean describes no player and measures nothing.}
\\
& \notesprompt{Now settle the hook vote --- and run the test on three more columns of digits.}
\begin{probgrid}{|Y|Y|}
\pcell{10}{\textbf{Vote again:} is the blog's conclusion correct? \; yes \; no. What type is Jersey \#, and what does the $14$ tell you about the team?}{1.7cm}{No. Jersey \# is categorical --- a label in digits --- so 14 describes no player.} & \pcell{11}{A school lists grade levels $9, 10, 11, 12$; the mean is $42/4 = 10.5$. Who is in grade $10.5$? What type is Grade level?}{1.4cm}{Nobody. Grade level is categorical --- a label in digits.} \\ \hline
\end{probgrid}
\\
& \begin{probgrid}*{|Y|Y|}
\pcell{12}{Two ZIP codes, $22030$ and $22201$, have a mean of $22115.5$. What place is that? What type is ZIP code?}{1.2cm}{No place at all --- a ZIP code identifies an area. Categorical.} & \pcell{13}{Height (in) is digits too: $420/6 = 70$. Who or what does $70$ inches describe? What type is Height?}{1.3cm}{A typical player's height, in inches. Quantitative.} \\ \hline
\end{probgrid}
\\ \hline
% ── ROW 5: GUIDED PRACTICE — we do, together, a fresh context ────────────────
\mainidea[Guided practice]{Student Life Survey} &
Six Riverbend students filled out the Student Life Survey. Theo says \emph{``Grade is
quantitative --- it's numbers.''} A \textbf{statistical question} is one whose answers
\emph{vary} from individual to individual, so it takes data to answer.
\\
& {\centering
\begin{tikzpicture}[scale=1.0]
  \fill[goldbg] (1.75,-3.15) rectangle (2.65,0);
  \draw[goldacc, thick] (1.75,-3.15) rectangle (2.65,0);
  \draw[thick] (0,-0.45) -- (12.2,-0.45);
  \draw[thick] (0,-3.15) -- (12.2,-3.15);
  \foreach \r/\pa/\pb/\pc/\pd/\pe/\pf in {
      0/\textbf{Student}/\textbf{Grade}/\textbf{Commute (min)}/\textbf{Lunch}/\textbf{Siblings}/\textbf{Club member},
      1/Alina/10/20/Home/2/Yes,
      2/Devon/11/8/Cafeteria/0/No,
      3/Marisol/9/35/Home/3/Yes,
      4/Theo/12/15/Cafeteria/1/Yes,
      5/Priya/10/5/Off campus/2/No,
      6/Jamal/11/25/Cafeteria/4/Yes} {
    \node[anchor=west, font=\footnotesize] at (0.1,{-0.45*\r-0.225}) {\pa};
    \node[font=\footnotesize] at (2.2,{-0.45*\r-0.225}) {\pb};
    \node[font=\footnotesize] at (4.4,{-0.45*\r-0.225}) {\pc};
    \node[anchor=west, font=\footnotesize] at (6.1,{-0.45*\r-0.225}) {\pd};
    \node[font=\footnotesize] at (9.0,{-0.45*\r-0.225}) {\pe};
    \node[font=\footnotesize] at (11.1,{-0.45*\r-0.225}) {\pf};
  }
\end{tikzpicture}\par}
\\
& \notesprompt{We work these together.}
\begin{probgrid}{|Y|Y|}
\pcell{14}{How many individuals does the table describe, and how many variables? Name one of each.}{1.4cm}{6 individuals; 5 variables --- e.g., Alina and Commute.} & \pcell{15}{Classify Commute, Lunch, and Grade. Then run the digit test on Grade ($63/6 = 10.5$) and fix Theo's reason.}{1.8cm}{Commute: quantitative. Lunch: categorical. Grade: categorical --- nobody is 10.5.} \\ \hline
\end{probgrid}
\\
& \begin{probgrid}*{|Y|Y|}
\pcell{16}{Find the mean commute, then say it as data --- who, what, units.}{1.4cm}{$108/6 = 18$ minutes, for the six surveyed.} & \pcell{17}{\emph{How long is Priya's commute?} \emph{How long are Riverbend students' commutes?} Which takes data? Rewrite the other.}{1.6cm}{The second --- answers vary. Rewrite: \emph{How long are Riverbend commutes?}} \\ \hline
\end{probgrid}
\\ \hline
\end{guidednotes}

% ── THE NOTES END AT GUIDED PRACTICE ─────────────────────────────────────────
% There is deliberately no "you do" here: ../homework/ is the individual practice,
% started in class in the last ten minutes while the teacher circulates.

\end{document}
